\documentclass[aspectratio=169,10pt]{beamer}

\mode<presentation>{
	\usetheme{Pittsburgh}
	\usecolortheme{beaver}
}

\usepackage{graphicx}
\usepackage{booktabs}
\usepackage{tabularx}
\usepackage{array}
\usepackage{multicol}
\usepackage{amsmath,bm}
\usepackage{fontspec}
\usepackage{xeCJK}
\usepackage{unicode-math}
\usepackage{ragged2e}
\usepackage{adjustbox}
\usepackage{xcolor}

\setsansfont{Calibri}[
	Path=../Font/,
	Scale=0.9,
	Extension=.ttf,
	UprightFont=*-Regular,
	BoldFont=*-Bold,
	ItalicFont=*-Italic,
	BoldItalicFont=*-BoldItalic
]
\setmathfont{Libertinus Math}
\setCJKmainfont{Songti SC}
\setCJKsansfont{Songti SC}

\justifying\let\raggedright\justifying
\setbeamertemplate{navigation symbols}{}
\setbeamersize{text margin left=8mm,text margin right=8mm}
\setbeamerfont{frametitle}{size=\large,series=\bfseries}
\setbeamerfont{framesubtitle}{size=\small}
\setbeamercolor{postit}{fg=black,bg=yellow!18}
\setbeamercolor{note}{fg=black,bg=gray!8}
\setbeamercolor{alerted text}{fg=red!70!black}

\newcolumntype{Y}{>{\raggedright\arraybackslash}X}
\newcommand{\papermax}{Nartea, Kong, and Wu (2017)}
\newcommand{\paperivol}{Gu, Kang, and Xu (2018)}
\newcommand{\tabnote}[1]{\vspace{2pt}{\scriptsize\textcolor{gray!70!black}{#1}}}
\newcommand{\tighttable}{\scriptsize\setlength{\tabcolsep}{3pt}\renewcommand{\arraystretch}{1.12}}
\newcommand{\tinytab}{\tiny\setlength{\tabcolsep}{2.2pt}\renewcommand{\arraystretch}{1.08}}

\title[China Anomaly]{\textbf{中国典型 anomaly 的自然研究文章}}
\subtitle{Extreme Returns, IVOL, and Limits of Arbitrage in China}
\author[Cheng Hang]{Cheng Hang}
\institute[DUFE]{School of Fintech\\Dongbei University of Finance \& Economics}
\date{\today}

\AtBeginSection[]
{
	\begin{frame}{Content}
		\small
		\tableofcontents[currentsection,sectionstyle=show/shaded,subsectionstyle=show/show/hide]
		\addtocounter{framenumber}{-1}
	\end{frame}
}

\begin{document}

\begin{frame}
	\titlepage
\end{frame}

\begin{frame}{为什么把这两篇放在一起读}
	这两篇文章的可读之处，不在于又报告了一个中国异象，而在于把中国市场当成研究经典 anomaly 机制的天然实验场：
	\begin{itemize}
		\item 中国股市散户成交占比高、投机色彩浓，便于观察彩票偏好如何进入横截面定价。
		\item 涨跌停、卖空限制、两融标的、股指期货覆盖不足等制度摩擦同时存在，便于观察套利限制如何放大 mispricing。
		\item 两篇文章背后都在追问同一个问题：经典 anomaly 究竟是风险补偿，是投资者偏好，还是价格没法被及时纠正？
		\item 怎么样讲好中国Story？
	\end{itemize}
	\vspace{4pt}
	\begin{beamercolorbox}[sep=0.45em,wd=\linewidth,rounded=true]{postit}
		评价这类文章好不好，要看\textbf{中国的制度细节是否真正参与了识别}，而不是把美国 anomaly 换个样本重跑一遍。
	\end{beamercolorbox}


\end{frame}

\begin{frame}{两篇文献的分工}
	\tighttable
	\begin{tabularx}{\textwidth}{lYY}
		\toprule
		文章 & 研究的 anomaly & 中国市场提供的识别条件 \\
		\midrule
		\papermax & 上月最大日收益 MAX 越高，下月收益越低 & 如果投资者偏好 lottery-like 股票，高 MAX 股票会被追捧并短期高估；中国散户主导的交易环境是天然检验场。\\
		\paperivol & 特质波动 IVOL 越高，未来收益越低 & 如果高 IVOL 代表被高估但难以套利，那么套利越受限，负 IVOL premium 应越强、越持久；中国的交易约束恰好提供了机制变量。\\
		\bottomrule
	\end{tabularx}
	\vspace{6pt}
	\begin{center}
		\small 可以这样定位：MAX 抓的是极端收益意义上的 lottery payoff；IVOL 抓的是难以套利的 mispricing。
	\end{center}
\end{frame}

\begin{frame}{怎么看后面的表}
	\begin{itemize}
		\item 后面每张表都对应三个问题：\textbf{这张表想检验什么？关键数字是多少？它在多大程度上支持文章主张？}
		\item 排序表不必照搬全部组合收益，盯住极端组的 spread 即可：MAX 表看 HMAX--LMAX，IVOL 表看 P5--P1 或 D10--D1。
		\item 技术上反复出现的是三件事：portfolio sorting、Fama--MacBeth 回归，以及风险调整后的 alpha。
	\end{itemize}
	\vspace{4pt}
	\begin{beamercolorbox}[sep=0.45em,wd=\linewidth,rounded=true]{note}
		读 anomaly 实证文章的顺序：先看极端组 spread 是不是稳定；再看加入控制变量后，spread 或回归斜率是否还显著。
	\end{beamercolorbox}
\end{frame}

\section{Nartea et al. (2017): MAX Effect}

\begin{frame}{文献一：先用一句话概括}
	\textbf{\papermax, Journal of Banking \& Finance, 2017}\\
	\textit{Do extreme returns matter in emerging markets? Evidence from the Chinese stock market}
	\vspace{6pt}
	\begin{itemize}
		\item 问题：美国和欧洲都报告了高 MAX 股票未来收益更低的现象，中国是否也存在？
		\item 思路：高 MAX 股票具有 lottery-like payoff，偏好彩票特征的投资者会把价格抬高，后续收益自然偏低。
		\item 结论：中国确实存在显著的负 MAX 效应；并且在中国，MAX 没有把 IVOL 效应吃掉，两者各自独立存在。
	\end{itemize}
\end{frame}

\begin{frame}{数据和样本}
	\begin{itemize}
		\item 数据：日频与月频收益、会计数据均来自 CSMAR。
		\item 样本：沪深 A 股，1997:01--2014:12；公司数由 515 增长到 2353，平均每月 1383 家，共 298{,}710 个 firm-month observations。
		\item 剔除：保留已退市股票，避免 survivorship bias。
		\item 处理：剔除月收益超过 200\% 的样本；IPO 首日收益不进入日频计算；日收益超出涨跌停实际范围的观测做删除或合并处理。
	\end{itemize}
	\tabnote{这些处理基本都是围绕涨跌停制度展开的，否则 MAX 会被价格限制机械截断，识别不出真正的极端收益。}
\end{frame}

\begin{frame}[c]{主要变量}
	\tighttable
	\begin{tabularx}{\textwidth}{lY}
		\toprule
		变量 & 定义 \\
		\midrule
		MAX & 上一自然月内最大单日收益；若连续触及涨跌停，则将连续涨跌停日的收益加总，近似无价格限制时的极端收益。\\
		Size & 月末市值的对数。\\
		BM & 第 $t-6$ 期的 book-to-market ratio。\\
		Momentum & $t-12$ 至 $t-2$ 的 11 个月累计收益。\\
		REV & 第 $t-1$ 期月收益，用来控制短期反转。\\
		SSKEW / ISKEW & 参照 Harvey and Siddique 的系统偏度与特质偏度。\\
		IV & 上月日收益对 FF3 回归后残差的标准差。\\
		ILLIQ & Amihud 非流动性指标。\\
		BETA & 对当期、滞后、领先市场收益回归得到的三个系数之和。\\
		\bottomrule
	\end{tabularx}
\end{frame}

\begin{frame}{实证设计}
	\begin{enumerate}
		\item 单变量排序：每月按 MAX 分五组，比较 High MAX 与 Low MAX 在下个月的收益以及 FF3 alpha。
		\item 双重排序：先按控制变量分三组，再在组内按 MAX 分五组，最后在 MAX 维度上做平均；控制变量包括 size、BM、momentum、REV、CP、skewness、IV、ILLIQ、beta。
		\item Fama--MacBeth 回归：用滞后一期变量预测当月收益，并依次给出二变量和多变量规格。
		\item 持有期延长：把持有期改为 3 个月、6 个月，看价格回归是否存在较长滞后。
	\end{enumerate}
	\vspace{4pt}
	\[
	R_{i,t} = \beta_{0,t-1}+\beta_{1,t-1}MAX_{i,t-1}+\cdots+\beta_{11,t-1}BETA_{i,t-1}+\varepsilon_{i,t-1}.
	\]
\end{frame}

\begin{frame}{主要结果：MAX 与 IVOL 在中国并存}
	\begin{itemize}
		\item High MAX 组合下月收益显著低于 Low MAX 组合，1 个月、6 个月持有期都成立。
		\item 加入大量控制变量后，MAX 的负斜率依然显著；仅当把 IV 单独放进回归时，MAX 在部分规格里会被削弱。
		\item 反过来，把 MAX 放进 IVOL 的回归里，IVOL 的负斜率也仍然存在，说明 MAX 不能简单替代 IVOL。
		\item 这一结论与 Bali et al. (2011) 的美国证据不同：在美国，MAX 被视为 ``真效应''、IVOL 更像它的代理；在中国，两者更像两个相互独立的异象。
	\end{itemize}
\end{frame}

\begin{frame}{Table 1：MAX 单变量组合}
	\tighttable
	\begin{tabular}{lrrrr}
		\toprule
		& EW return & EW FF3 alpha & VW return & VW FF3 alpha \\
		\midrule
		Low MAX & 0.0183 & 0.0029 & 0.0134 & 0.0024 \\
		High MAX & 0.0072 & -0.0085 & 0.0074 & -0.0045 \\
		High--Low & -0.0111 & -0.0114 & -0.0060 & -0.0070 \\
		t-stat & -6.21 & -6.26 & -2.30 & -2.57 \\
		\bottomrule
	\end{tabular}
	\vspace{6pt}
	\begin{itemize}
		\item 高 MAX 组下月跑输显著，扣除 FF3 之后仍然成立。
		\item 等权 spread 比市值加权更大，提示这一异象在小股票中更强；不过市值加权仍显著，所以并不是纯粹的 microcap 现象。
	\end{itemize}
\end{frame}

\begin{frame}{Table 2：高 MAX 股票长什么样}
	\tinytab
	\begin{tabular}{lrrrrrrrrrr}
		\toprule
		& Size & BM & Mom & REV & CP & SSKEW & ISKEW & IV & ILLIQ & BETA\\
		\midrule
		Low MAX & 14.313 & 0.443 & 0.180 & -0.032 & 10.042 & -1.517 & 0.117 & 0.012 & 0.0026 & 0.937\\
		High MAX & 14.175 & 0.391 & 0.284 & 0.079 & 11.193 & -1.961 & 0.925 & 0.025 & 0.0025 & 1.040\\
		High--Low & -0.138 & -0.052 & 0.104 & 0.111 & 1.151 & -0.444 & 0.808 & 0.013 & -0.0002 & 0.103\\
		t-stat & -19.31 & -38.94 & 27.77 & 113.19 & 32.97 & -12.98 & 155.46 & 269.69 & -6.27 & 58.61\\
		\bottomrule
	\end{tabular}
	\vspace{5pt}
	\begin{itemize}
		\item 高 MAX 股票市值更小、BM 更低、过去收益更强、IV 更高、特质偏度更高、beta 也更高。
		\item 这个画像非常像典型的 lottery-like 股票；也提醒我们：后面的回归必须同时控制 IV 和偏度。
	\end{itemize}
\end{frame}

\begin{frame}{Table 3：双重排序控制后 MAX alpha}
	\tinytab
	\begin{tabular}{lrrrr}
		\toprule
		控制变量 & EW HMAX--LMAX & t & VW HMAX--LMAX & t \\
		\midrule
		Size & -0.0133 & -7.18 & -0.0120 & -6.09\\
		BM & -0.0160 & -8.34 & -0.0124 & -4.92\\
		Momentum & -0.0129 & -6.61 & -0.0096 & -3.93\\
		REV & -0.0139 & -6.46 & -0.0094 & -3.56\\
		CP & -0.0125 & -6.87 & -0.0093 & -3.64\\
		SSKEW & -0.0135 & -6.39 & -0.0095 & -3.95\\
		ISKEW & -0.0105 & -4.29 & -0.0048 & -1.49\\
		IV & -0.0028 & -0.92 & -0.0013 & -0.37\\
		ILLIQ & -0.0126 & -6.85 & -0.0096 & -4.33\\
		BETA & -0.0124 & -5.87 & -0.0077 & -2.81\\
		\bottomrule
	\end{tabular}
	\vspace{4pt}
	\begin{itemize}
		\item 除了控制 IV、以及部分 ISKEW 规格之外，MAX 的 spread 都显著为负。
		\item 真正的挑战来自 IV：高 MAX 与高 IV 高度相关，所以在后面的回归里必须同时放入两者。
	\end{itemize}
\end{frame}

\begin{frame}{Table 4：单变量 Fama--MacBeth}
	\tighttable
	\begin{tabular}{lrrrrrrrrrrr}
		\toprule
		& MAX & Size & BM & Mom & REV & CP & SSKEW & ISKEW & IV & ILLIQ & BETA\\
		\midrule
		Coef. & -0.155 & -0.002 & 0.012 & 0.000 & -0.049 & 0.000 & 0.001 & -0.004 & -0.948 & 10.068 & -0.003\\
		t & -7.21 & -1.37 & 2.33 & 0.06 & -4.64 & -1.19 & 1.96 & -4.03 & -9.22 & 2.01 & -0.58\\
		\bottomrule
	\end{tabular}
	\vspace{6pt}
	\begin{itemize}
		\item 单独放进回归时，MAX 的系数显著为负；REV、ISKEW、IV 同样为负，BM 和 ILLIQ 为正。
		\item 这张表只告诉我们每个变量\emph{单独}能解释什么，还不能作为机制判断的依据。
	\end{itemize}
\end{frame}

\begin{frame}{Table 5：MAX 加控制变量后的 Fama--MacBeth}
	\tinytab
	\begin{tabular}{lrr}
		\toprule
		规格 & MAX coefficient & t-stat \\
		\midrule
		MAX + Size & -0.161 & -8.42\\
		MAX + BM & -0.157 & -8.34\\
		MAX + Momentum & -0.151 & -7.70\\
		MAX + REV & -0.145 & -5.76\\
		MAX + CP & -0.161 & -8.17\\
		MAX + SSKEW & -0.158 & -8.01\\
		MAX + ISKEW & -0.130 & -5.50\\
		MAX + IV & 0.004 & 0.12\\
		MAX + ILLIQ & -0.149 & -7.02\\
		MAX + BETA & -0.155 & -7.10\\
		All controls & -0.072 & -3.72\\
		\bottomrule
	\end{tabular}
	\hfill
	\begin{minipage}{0.44\textwidth}
		\small
		\begin{itemize}
			\item 单独和 IV 放在一起时，MAX 失去显著性。
			\item 把所有变量一起放进去之后，MAX 重新显著为负，说明它不能被常规特征完全吸收。
			\item 同一规格里 IV = -0.182，t = -3.07，可见 MAX 和 IV 是并存关系，而不是互相替代。
		\end{itemize}
	\end{minipage}
\end{frame}

\begin{frame}{Table 6：6 个月持有期的 MAX 组合}
	\tighttable
	\begin{tabular}{lrrrr}
		\toprule
		& EW return & EW FF3 alpha & VW return & VW FF3 alpha\\
		\midrule
		Low MAX & 0.0160 & 0.0020 & 0.0122 & 0.0028\\
		High MAX & 0.0102 & -0.0048 & 0.0069 & -0.0040\\
		High--Low & -0.0058 & -0.0067 & -0.0052 & -0.0068\\
		t-stat & -4.72 & -5.85 & -2.93 & -4.43\\
		\bottomrule
	\end{tabular}
	\vspace{6pt}
	\begin{itemize}
		\item 负 MAX 效应延续到 6 个月，提示价格向基本面收敛存在较长的滞后。
		\item 中国卖空受限的现实，可能让高估状态的纠偏过程更慢。
	\end{itemize}
\end{frame}

\begin{frame}{Table 7：6 个月持有期双重排序}
	\tinytab
	\begin{tabular}{lrrrr}
		\toprule
		控制变量 & EW HMAX--LMAX & t & VW HMAX--LMAX & t\\
		\midrule
		Size & -0.0062 & -5.85 & -0.0061 & -5.43\\
		BM & -0.0081 & -7.41 & -0.0084 & -5.57\\
		Momentum & -0.0059 & -5.71 & -0.0060 & -4.32\\
		REV & -0.0075 & -6.11 & -0.0065 & -4.22\\
		CP & -0.0061 & -6.16 & -0.0062 & -4.76\\
		SSKEW & -0.0059 & -5.39 & -0.0057 & -4.01\\
		ISKEW & -0.0052 & -3.80 & -0.0045 & -2.55\\
		IV & -0.0018 & -0.84 & -0.0006 & -0.25\\
		ILLIQ & -0.0058 & -5.51 & -0.0048 & -4.11\\
		BETA & -0.0058 & -5.23 & -0.0046 & -3.34\\
		\bottomrule
	\end{tabular}
	\vspace{4pt}
	\begin{itemize}
		\item 控制其他变量后，6 个月持有期下的 MAX spread 同样显著为负。
		\item 再次出现的现象是：只有控制 IV 时效应才被削弱。这就把 MAX 与 IVOL 的关系拉成全文的关键张力。
	\end{itemize}
\end{frame}

\begin{frame}{Table 8--9：6 个月 Fama--MacBeth}
	\begin{columns}[t]
		\begin{column}{0.48\textwidth}
			\textbf{Table 8：单变量}
			\scriptsize
			\begin{tabular}{lrr}
				\toprule
				变量 & Coef. & t\\
				\midrule
				MAX & -0.581 & -8.01\\
				Size & -0.020 & -2.56\\
				BM & 0.049 & 2.47\\
				REV & -0.100 & -3.17\\
				ISKEW & -0.008 & -4.27\\
				IV & -2.539 & -7.12\\
				ILLIQ & 60.299 & 4.09\\
				\bottomrule
			\end{tabular}
		\end{column}
		\begin{column}{0.52\textwidth}
			\textbf{Table 9：MAX 加控制变量}
			\scriptsize
			\begin{tabular}{lrr}
				\toprule
				规格 & MAX coefficient & t\\
				\midrule
				+ Size & -0.556 & -9.09\\
				+ BM & -0.535 & -8.68\\
				+ Momentum & -0.523 & -8.73\\
				+ REV & -0.587 & -8.39\\
				+ CP & -0.514 & -8.79\\
				+ SSKEW & -0.554 & -7.93\\
				+ ISKEW & -0.561 & -6.41\\
				+ IV & -0.292 & -3.76\\
				All controls & -0.155 & -2.16\\
				\bottomrule
			\end{tabular}
		\end{column}
	\end{columns}
	\par\vspace{5pt}
	\small 6 个月持有期下的回归进一步印证了主结论：MAX 和 IV 都保留独立的解释力。
\end{frame}

\begin{frame}{Table 10：规模组和子样本}
	\tinytab
	\begin{tabular}{lrrrrr}
		\toprule
		变量 & Small & Medium & Big & 1997--2004 & 2005--2014\\
		\midrule
		MAX & -0.094 & -0.111 & -0.270*** & -0.092 & -0.176**\\
		Size & -0.028** & -0.026*** & -0.014** & -0.017 & -0.017**\\
		BM & 0.069*** & 0.035** & 0.042** & 0.097*** & 0.028*\\
		Momentum & -0.009 & 0.001 & 0.017 & 0.084*** & -0.023\\
		REV & -0.007 & 0.014 & 0.035 & 0.091*** & -0.014\\
		CP & -0.007** & -0.002 & 0.000 & -0.001 & -0.001\\
		SSKEW & 0.004*** & 0.004*** & 0.007*** & 0.004** & 0.005**\\
		ISKEW & -0.003 & 0.000 & -0.002 & -0.006*** & -0.001\\
		IV & -1.712*** & -2.033*** & -1.055*** & -2.299*** & -1.387***\\
		ILLIQ & 28.981*** & 34.390*** & 35.649** & 1.375 & 50.158***\\
		BETA & 0.014 & -0.003 & -0.013 & 0.005 & -0.012\\
		\bottomrule
	\end{tabular}
	\vspace{4pt}
	\begin{itemize}
		\item MAX 在大市值组和 2005--2014 子样本中更显著。
		\item 作者的解释是：A 股中大公司往往股价绝对水平较低，更容易成为散户彩票式交易的对象。
	\end{itemize}
\end{frame}

\begin{frame}{附录 Table 1A--1B：FF3 因子的描述统计}
	\begin{columns}[t]
		\begin{column}{0.52\textwidth}
			\textbf{Table 1A：FF3 因子描述统计}
			\scriptsize
			\begin{tabular}{lrrr}
				\toprule
				& Market & SMB & HML\\
				\midrule
				Average return & 0.93\% & 0.82\% & 0.19\%\\
				t-stat & 1.59 & 2.98 & 0.95\\
				Minimum & -26.85\% & -17.30\% & -11.81\%\\
				Median & 0.73\% & 1.12\% & 0.21\%\\
				Maximum & 36.29\% & 10.08\% & 14.84\%\\
				Skewness & 0.33 & -0.55 & -0.11\\
				Kurtosis & 4.67 & 4.51 & 7.27\\
				\bottomrule
			\end{tabular}
		\end{column}
		\begin{column}{0.45\textwidth}
			\textbf{Table 1B：FF3 因子相关矩阵}
			\scriptsize
			\begin{tabular}{lrrr}
				\toprule
				& Market & SMB & HML\\
				\midrule
				Market & 1.00 & 0.07 & 0.03\\
				SMB & 0.19 & 1.00 & -0.24\\
				HML & 0.04 & -0.25 & 1.00\\
				\bottomrule
			\end{tabular}
			\par\vspace{4pt}
			\scriptsize 上三角为 Pearson，下三角为 Spearman。
		\end{column}
	\end{columns}
\end{frame}

\section{Gu et al. (2018): IVOL and Limits of Arbitrage}

\begin{frame}{文献二：先用一句话概括}
	\textbf{\paperivol, Journal of Banking \& Finance, 2018}\\
	\textit{Limits of arbitrage and idiosyncratic volatility: Evidence from China stock market}
	\vspace{6pt}
	\begin{itemize}
		\item 问题：为什么高 IVOL 股票未来收益更低？中国市场的套利限制会不会把这一负 premium 进一步放大？
		\item 思路：构造一个综合的 limits-of-arbitrage 指数，把中国特有的交易约束与一般市场摩擦放在同一指标内。
		\item 结论：在套利受限更严重的股票中，负 IVOL premium 更强、也更持久；短期反转、彩票偏好和套利不对称都不足以单独解释这一现象。
	\end{itemize}
\end{frame}

\begin{frame}{样本和变量}
	\begin{itemize}
		\item 样本：沪深 A 股，2002:01--2012:12；收益与会计数据来自 CSMAR，FF3 因子来自 RESSET。
		\item 剔除：金融行业公司和 ST 公司。
		\item 会计数据匹配：第 $t-1$ 财年的会计变量，匹配到 $t$ 年 7 月至 $t+1$ 年 6 月的收益区间。
		\item IVOL：在每个月内，用日度超额收益对 FF3 因子做回归，取残差标准差。
		\item MAX5：上一月日收益最高的 5 天的平均收益，作为 lottery-like payoff 的代理。
	\end{itemize}
	\[
	R_{i,d}-r_{f,d}=\alpha_i+\beta_i MKT_d+s_i SMB_d+h_i HML_d+\varepsilon_{i,d}.
	\]
\end{frame}

\begin{frame}{套利限制指数：六个组件}
	\tighttable
	\begin{tabularx}{\textwidth}{lY}
		\toprule
		组件 & 经济含义 \\
		\midrule
		Price-limit-hitting & 当月是否触及过涨跌停，触及越频繁，价格发现越受阻。\\
		MTSS availability & 是否不在融资融券标的之内；不能卖空，套利者就难以做空被高估的股票。\\
		CSI 300 membership & 是否不属于沪深 300 成分股；非成分股更难用股指期货对冲。\\
		Amihud illiquidity & 非流动性是否高于当月中位数。\\
		Trading volume & 交易量是否低于或等于中位数。\\
		Analyst coverage & 分析师覆盖是否低于或等于中位数。\\
		\bottomrule
	\end{tabularx}
	\par\vspace{5pt}
	\small 每个组件识别到高套利限制时取 1；综合指数为可用组件的平均值，指数越大，套利限制越严重。
\end{frame}

\begin{frame}{Table 1：描述统计和相关性}
	\begin{columns}[t]
		\begin{column}{0.48\textwidth}
			\textbf{Panel A：描述统计}
			\scriptsize
			\begin{tabular}{lrrrrr}
				\toprule
				& Mean & SD & Q1 & Median & Q3\\
				\midrule
				RET$_{t+1}$ & 1.34 & 9.26 & -4.69 & 0.22 & 6.18\\
				IVOL$_t$ & 1.79 & 0.71 & 1.28 & 1.67 & 2.18\\
				lnBM$_t$ & -1.27 & 0.96 & -1.40 & -1.15 & -0.69\\
				lnMV$_t$ & 0.08 & 0.54 & -0.39 & 0.02 & 0.80\\
				TURN$_t$ & 2.79 & 1.68 & 1.69 & 2.43 & 3.46\\
				MAX5$_t$ & 3.39 & 1.20 & 2.58 & 3.19 & 4.00\\
				\bottomrule
			\end{tabular}
		\end{column}
		\begin{column}{0.52\textwidth}
			\textbf{Panel B：关键相关性}
			\scriptsize
			\begin{itemize}
				\item RET$_{t+1}$ 与 IVOL：Pearson -0.07，Spearman -0.10。
				\item RET$_{t+1}$ 与 MAX5：Pearson -0.06，Spearman -0.08。
				\item IVOL 与 MAX5：Pearson 0.76，Spearman 0.72。
				\item TURN 与 IVOL：0.29；TURN 与 MAX5：0.30。
			\end{itemize}
			\par\vspace{4pt}
			\small 一个重要信息：IVOL 与 MAX5 高度相关，后面任何关于 IVOL 的检验都必须把 MAX5 一并控制。
		\end{column}
	\end{columns}
\end{frame}

\begin{frame}{Table 2：单变量 IVOL 组合}
	\tighttable
	\begin{tabular}{lrrrr}
		\toprule
		& VW raw & VW FF3 adj. & EW raw & EW FF3 adj.\\
		\midrule
		Quintile P5--P1 & -1.02 & -1.09 & -1.79 & -1.76\\
		t-stat & -2.60 & -3.14 & -7.63 & -8.05\\
		Decile D10--D1 & -1.32 & -1.40 & -2.21 & -2.20\\
		t-stat & -2.77 & -3.37 & -7.51 & -7.84\\
		\bottomrule
	\end{tabular}
	\vspace{6pt}
	\begin{itemize}
		\item 高 IVOL 股票下月收益显著低于低 IVOL 股票。
		\item 在中国，IVOL premium 的经济量级相当可观：等权风险调整后的 D10--D1 已经达到每月 -2.20\%。
	\end{itemize}
\end{frame}

\begin{frame}{Table 3：套利限制越强，IVOL premium 越强}
	\tinytab
	\begin{tabular}{lrrrrrrrr}
		\toprule
		& \multicolumn{2}{c}{VW raw} & \multicolumn{2}{c}{VW FF3} & \multicolumn{2}{c}{EW raw} & \multicolumn{2}{c}{EW FF3}\\
		& P5--P1 & t & P5--P1 & t & P5--P1 & t & P5--P1 & t\\
		\midrule
		Low LA & -0.43 & -0.99 & -0.47 & -1.19 & -0.99 & -3.15 & -0.95 & -3.40\\
		Medium LA & -1.65 & -4.70 & -1.53 & -4.66 & -2.09 & -7.62 & -2.01 & -7.66\\
		High LA & -2.35 & -8.52 & -2.18 & -8.42 & -2.40 & -9.33 & -2.29 & -9.13\\
		High--Low & -1.92 & -4.31 & -1.71 & -4.03 & -1.41 & -4.34 & -1.34 & -4.25\\
		\bottomrule
	\end{tabular}
	\vspace{5pt}
	\begin{itemize}
		\item 这是全文的关键表：在低套利限制组里，市值加权 IVOL premium 几乎不显著；进入高套利限制组后，则强烈显著。
		\item 跨高低套利限制组的 spread 差异（High--Low）同样显著为负，说明套利限制确实在放大 IVOL anomaly。
	\end{itemize}
\end{frame}

\begin{frame}{Table 4：三组件指数的稳健性}
	\tinytab
	\begin{tabular}{lrrrr}
		\toprule
		& \multicolumn{2}{c}{Panel A: price limit + CSI + illiquidity} & \multicolumn{2}{c}{Panel B: MTSS + volume + coverage}\\
		& VW FF3 P5--P1 & EW FF3 P5--P1 & VW FF3 P5--P1 & EW FF3 P5--P1\\
		\midrule
		Low LA & -0.75 & -1.32 & -0.69 & -1.00\\
		Medium LA & -1.14 & -1.98 & -1.45 & -1.90\\
		High LA & -1.98 & -2.32 & -2.01 & -2.18\\
		High--Low & -1.23 & -1.00 & -1.31 & -1.18\\
		t-stat High--Low & -3.09 & -3.46 & -3.16 & -3.62\\
		\bottomrule
	\end{tabular}
	\par\vspace{5pt}
	\small 无论换用哪三个组件来构造 LA 指数，高套利限制组的负 IVOL premium 都会更强，结论稳健。
\end{frame}

\begin{frame}{Table 5：去掉任一组件后的稳健性}
	\tinytab
	\begin{tabular}{lrrrrrr}
		\toprule
		去掉的组件 & Price limit & MTSS & CSI & Illiq. & Volume & Coverage\\
		\midrule
		VW raw High--Low & -1.31 & -1.92 & -1.92 & -1.96 & -1.62 & -1.91\\
		t & -3.02 & -4.32 & -4.29 & -4.69 & -3.99 & -4.63\\
		VW FF3 High--Low & -1.17 & -1.71 & -1.71 & -1.78 & -1.44 & -1.66\\
		t & -2.79 & -4.03 & -4.03 & -4.44 & -3.74 & -4.31\\
		EW raw High--Low & -0.89 & -1.42 & -1.42 & -1.71 & -1.15 & -1.41\\
		t & -2.77 & -4.35 & -4.35 & -5.50 & -3.94 & -4.40\\
		EW FF3 High--Low & -0.86 & -1.35 & -1.35 & -1.64 & -1.11 & -1.29\\
		t & -2.74 & -4.27 & -4.28 & -5.33 & -3.86 & -4.13\\
		\bottomrule
	\end{tabular}
	\par\vspace{5pt}
	\small 任意去掉一个组件后，结果都不会被打掉；综合 LA 指数的结论并非由某一个变量单独驱动。
\end{frame}

\begin{frame}{Table 6：IVOL premium 的持续性}
	\tinytab
	\begin{tabular}{lrrrrrr}
		\toprule
		& t+1 & t+2 & t+3 & t+4 & t+5 & t+6\\
		\midrule
		\multicolumn{7}{l}{Raw return P5--P1}\\
		Low LA & -0.43 & -0.67 & -0.88 & -1.26 & -1.16 & -1.03\\
		Medium LA & -1.65 & -2.54 & -2.83 & -3.73 & -3.99 & -4.36\\
		High LA & -2.35 & -3.37 & -3.96 & -4.40 & -4.63 & -4.96\\
		High--Low & -1.92 & -2.69 & -3.08 & -3.14 & -3.47 & -3.93\\
		\midrule
		\multicolumn{7}{l}{FF3 adjusted P5--P1}\\
		Low LA & -0.47 & -0.60 & -0.64 & -0.71 & -0.64 & -0.85\\
		Medium LA & -1.53 & -2.22 & -2.27 & -2.93 & -3.01 & -3.53\\
		High LA & -2.18 & -2.98 & -3.46 & -3.98 & -3.93 & -4.19\\
		High--Low & -1.71 & -2.38 & -2.82 & -3.27 & -3.29 & -3.34\\
		\bottomrule
	\end{tabular}
	\par\vspace{5pt}
	\small 高套利限制组的负 IVOL premium 不仅更大，还在之后 6 个月里继续累积，没有出现明显反转。
\end{frame}

\begin{frame}{Table 7：Fama--MacBeth 控制套利限制}
	\tinytab
	\begin{tabular}{lrrr}
		\toprule
		变量 & (1) & (2) & (3)\\
		\midrule
		IVOL$^{ran}$ & -1.08*** & & -0.69\\
		& (-2.87) & & (-1.65)\\
		IVOL$^{ran}\times$Low LA & & -0.69 & \\
		& & (-1.65) & \\
		IVOL$^{ran}\times$Medium LA & & -1.57*** & -0.88**\\
		& & (-3.55) & (-2.29)\\
		IVOL$^{ran}\times$High LA & & -2.07*** & -1.37***\\
		& & (-4.90) & (-2.82)\\
		XRET$^{ran}$ & -0.99** & -0.96** & -0.96**\\
		lnMV$^{ran}$ & -0.88 & -0.84 & -0.84\\
		lnBM$^{ran}$ & 0.36 & 0.37 & 0.37\\
		TURN$^{ran}$ & -0.70* & -0.72* & -0.72*\\
		MAX5$^{ran}$ & 0.17 & 0.12 & 0.12\\
		$R^2$ & 14.9\% & 15.3\% & 15.3\%\\
		\bottomrule
	\end{tabular}
	\par\vspace{4pt}
	\small Fama--MacBeth 回归在系数层面重复了 Table 3 的结论：IVOL 的负系数在高 LA 组最强，并且在控制 MAX5 后依然成立。
\end{frame}

\begin{frame}{Table 8：短期反转不能解释 IVOL premium}
	\begin{columns}[t]
		\begin{column}{0.48\textwidth}
			\textbf{Panel A：时间序列 alpha}
			\scriptsize
			\begin{tabular}{lrrrr}
				\toprule
				& FF3 & +UMD & +STREV & +UMD+STREV\\
				\midrule
				Constant & -1.09*** & -1.19*** & -0.87** & -1.05***\\
				t & -2.82 & -3.32 & -2.13 & -2.78\\
				$R^2$ & 23.4\% & 28.7\% & 25.3\% & 29.0\%\\
				\bottomrule
			\end{tabular}
		\end{column}
		\begin{column}{0.52\textwidth}
			\textbf{Panel B：五因子调整后 P5--P1}
			\scriptsize
			\begin{tabular}{lrrrr}
				\toprule
				& VW & t & EW & t\\
				\midrule
				Low LA & -0.41 & -1.00 & -0.98 & -3.42\\
				Medium LA & -1.58 & -4.48 & -2.08 & -7.49\\
				High LA & -2.28 & -8.24 & -2.29 & -8.57\\
				High--Low & -1.87 & -4.34 & -1.31 & -4.13\\
				\bottomrule
			\end{tabular}
		\end{column}
	\end{columns}
	\par\vspace{5pt}
	\small 即便在 FF3 之外再加入动量和短期反转因子，主结论依然没有消失。
\end{frame}

\begin{frame}{Table 9：IVOL 与 MAX5 谁主导？}
	\tighttable
	\begin{tabular}{lrrrr}
		\toprule
		& EW raw & EW FF3 & VW raw & VW FF3\\
		\midrule
		\multicolumn{5}{l}{Panel A: sorting by IVOL}\\
		One-way D10--D1 & -2.21 & -2.20 & -1.32 & -1.40\\
		After controlling MAX5 & -1.44 & -1.35 & -0.97 & -0.90\\
		t-stat controlled & -4.77 & -4.70 & -2.92 & -3.00\\
		\midrule
		\multicolumn{5}{l}{Panel B: sorting by MAX5}\\
		One-way D10--D1 & -1.59 & -1.70 & -0.95 & -1.16\\
		After controlling IVOL & 0.16 & 0.03 & 0.17 & 0.00\\
		t-stat controlled & 0.46 & 0.08 & 0.41 & -0.01\\
		\bottomrule
	\end{tabular}
	\par\vspace{5pt}
	\small 控制 MAX5 后，IVOL 仍然显著；反过来控制 IVOL 后，MAX5 几乎完全消失。Gu et al. 据此认为 IVOL 才是更基础的定价变量。
\end{frame}

\begin{frame}{Table 10：控制 MAX5 后，LA--IVOL 仍成立}
	\tinytab
	\begin{tabular}{lrrrr}
		\toprule
		& VW raw & VW FF3 & EW raw & EW FF3\\
		\midrule
		\multicolumn{5}{l}{Panel A: LA and IVOL after controlling MAX5, P5--P1}\\
		Low LA & -0.46 & -0.36 & -0.89 & -0.75\\
		Medium LA & -1.22 & -1.10 & -1.40 & -1.32\\
		High LA & -1.57 & -1.43 & -1.47 & -1.37\\
		High--Low & -1.11 & -1.07 & -0.58 & -0.62\\
		\midrule
		\multicolumn{5}{l}{Panel B: LA and MAX5 after controlling IVOL, P5--P1}\\
		Low LA & -0.12 & -0.29 & 0.12 & -0.11\\
		Medium LA & -0.16 & -0.29 & -0.11 & -0.22\\
		High LA & -0.34 & -0.42 & -0.45 & -0.51\\
		High--Low & -0.22 & -0.13 & -0.57 & -0.40\\
		\bottomrule
	\end{tabular}
	\par\vspace{5pt}
	\small 三重排序进一步说明：lottery-like 收益特征并不能吸收掉 IVOL 与套利限制的交互效应。
\end{frame}

\begin{frame}{Table 11：套利不对称解释的边界}
	\tighttable
	\begin{tabular}{lrrrrrr}
		\toprule
		& \multicolumn{2}{c}{Low LA} & \multicolumn{2}{c}{Medium LA} & \multicolumn{2}{c}{High LA}\\
		& Under & Over & Under & Over & Under & Over\\
		\midrule
		Raw P5--P1 & -0.26 & -1.29 & -1.58 & -1.88 & -2.09 & -2.48\\
		t & -0.58 & -2.75 & -3.58 & -4.70 & -6.14 & -8.43\\
		Over--Under & \multicolumn{2}{c}{-1.04} & \multicolumn{2}{c}{-0.30} & \multicolumn{2}{c}{-0.39}\\
		t & \multicolumn{2}{c}{-2.20} & \multicolumn{2}{c}{-0.62} & \multicolumn{2}{c}{-1.21}\\
		\midrule
		FF3 P5--P1 & -0.26 & -1.29 & -1.45 & -1.75 & -1.91 & -2.34\\
		t & -0.64 & -2.75 & -3.53 & -4.55 & -5.77 & -8.29\\
		Over--Under & \multicolumn{2}{c}{-1.04} & \multicolumn{2}{c}{-0.30} & \multicolumn{2}{c}{-0.42}\\
		\bottomrule
	\end{tabular}
	\par\vspace{5pt}
	\small 在低 LA 组中，结果符合 Stambaugh et al.\ 的套利不对称：负 IVOL premium 主要集中在 overpriced 股票里。但是在高 LA 组中，无论 under- 还是 overpriced 股票，负 IVOL premium 都很显著——单靠套利不对称解释不充分。
\end{frame}

\begin{frame}{Table 12：回归检验套利不对称}
	\tinytab
	\begin{tabular}{lrr}
		\toprule
		变量 & (1) & (2)\\
		\midrule
		IVOL$^{ran}\times$Underpriced & -0.87* & \\
		& (-1.96) & \\
		IVOL$^{ran}\times$Overpriced & -1.60*** & \\
		& (-3.89) & \\
		IVOL$^{ran}\times$Low LA$\times$Under & & -0.50\\
		IVOL$^{ran}\times$Medium LA$\times$Under & & -1.61***\\
		IVOL$^{ran}\times$High LA$\times$Under & & -2.09***\\
		IVOL$^{ran}\times$Low LA$\times$Over & & -1.47***\\
		IVOL$^{ran}\times$Medium LA$\times$Over & & -1.76***\\
		IVOL$^{ran}\times$High LA$\times$Over & & -2.14***\\
		XRET$^{ran}$ & -0.99** & -0.96**\\
		TURN$^{ran}$ & -0.69* & -0.69*\\
		MAX5$^{ran}$ & 0.20 & 0.18\\
		$R^2$ & 15.3\% & 16.0\%\\
		\bottomrule
	\end{tabular}
	\par\vspace{5pt}
	\small 在高 LA 组中，无论 underpriced 还是 overpriced 子样本，IVOL 的交互项都显著为负，支持 ``高 LA 股票整体被高估'' 的解释。
\end{frame}

\begin{frame}{附录 A1--A3：LA 指数的底层质量}
	\begin{columns}[t]
		\begin{column}{0.45\textwidth}
			\textbf{A1：六个组件定义}
			\scriptsize
			\begin{itemize}
				\item LALIMIT：当月至少一次触及涨跌停。
				\item LAMTSS：未纳入融资融券卖空名单。
				\item LACSI：不属于沪深 300。
				\item LAAMIHUD：Amihud 非流动性高于中位数。
				\item LAVOL：交易量低于或等于中位数。
				\item LACOV：分析师覆盖低于或等于中位数。
			\end{itemize}
		\end{column}
		\begin{column}{0.55\textwidth}
			\textbf{A2--A3：分布和相关性}
			\scriptsize
			\begin{tabular}{lrr}
				\toprule
				LA index bin & Number & Percentage\\
				\midrule
				$[0,0.2)$ & 315 & 20.8\\
				$[0.2,0.4)$ & 275 & 18.2\\
				$[0.4,0.6)$ & 243 & 16.0\\
				$[0.6,0.8)$ & 408 & 27.0\\
				$[0.8,1]$ & 271 & 17.9\\
				\bottomrule
			\end{tabular}
			\par\vspace{4pt}
			\scriptsize A3 显示六个组件之间的相关性都显著；其中最高是 LAMTSS--LACSI = 0.76，LAAMIHUD--LAVOL = 0.68——这正是作者必须做替代指数稳健性检验的原因。
		\end{column}
	\end{columns}
\end{frame}

\begin{frame}{附录 A4：两组件指数下的稳健性}
	\tinytab
	\begin{tabular}{lrrrr}
		\toprule
		Panel & VW FF3 Low & VW FF3 High & VW FF3 High--Low & t\\
		\midrule
		A: price limit + volume & -0.55 & -1.90 & -1.35 & -3.26\\
		B: MTSS + analyst coverage & -0.53 & -2.42 & -1.89 & -4.03\\
		C: CSI + illiquidity & -0.98 & -2.00 & -1.02 & -2.45\\
		\bottomrule
	\end{tabular}
	\vspace{6pt}
	\tinytab
	\begin{tabular}{lrrrr}
		\toprule
		Panel & EW FF3 Low & EW FF3 High & EW FF3 High--Low & t\\
		\midrule
		A: price limit + volume & -1.16 & -2.37 & -1.21 & -4.41\\
		B: MTSS + analyst coverage & -1.04 & -2.76 & -1.73 & -4.83\\
		C: CSI + illiquidity & -1.47 & -2.02 & -0.54 & -1.96\\
		\bottomrule
	\end{tabular}
	\par\vspace{5pt}
	\small 即便把 LA 指数压缩到只有两个组件，``高 LA 组的 IVOL premium 更强'' 这一结论仍然成立。
\end{frame}

\begin{frame}{附录 A5--A6：控制 MAX5 与 overpricing 指标的构造}
	\begin{columns}[t]
		\begin{column}{0.48\textwidth}
			\textbf{A5：控制 MAX5 的 Fama--MacBeth}
			\scriptsize
			\begin{tabular}{lrrrrrr}
				\toprule
				& (1)&(2)&(3)&(4)&(5)&(6)\\
				\midrule
				IVOL & -0.90**&&-1.07**&-0.98***&&-1.08***\\
				MAX5 &&-0.75*&0.05&&-0.56&0.18\\
				XRET &&&&-0.89**&-0.92*&-0.98**\\
				TURN &&&&-0.62&-0.82**&-0.72*\\
				$R^2$ &3.1&3.4&5.3&13.4&14.0&14.9\\
				\bottomrule
			\end{tabular}
		\end{column}
		\begin{column}{0.52\textwidth}
			\textbf{A6：overpricing/underpricing 六个 anomaly}
			\scriptsize
			\begin{itemize}
				\item Accruals：应计项越高，未来收益越低。
				\item Net operating assets：净经营资产越高，未来收益越低。
				\item Gross profitability：盈利能力高，未来收益高。
				\item ROE：ROE 高，未来收益高。
				\item Momentum：6/1/6 动量。
				\item Composite equity issuance：股权发行越多，未来收益越低。
			\end{itemize}
		\end{column}
	\end{columns}
\end{frame}

\begin{frame}{附录 A7：高 LA 股票的估值更贵}
	\tighttable
	\begin{tabular}{lrrrrrrr}
		\toprule
		& Full & Low U & Low O & Med U & Med O & High U & High O\\
		\midrule
		\multicolumn{8}{l}{Panel A: median P/E}\\
		2002--2012 & 50.7 & 26.7 & 43.0 & 40.2 & 61.6 & 50.9 & 82.1\\
		2002--2007 & 50.4 & 24.2 & 40.0 & 34.6 & 69.4 & 49.8 & 92.2\\
		2008--2012 & 50.9 & 32.0 & 38.6 & 44.2 & 55.6 & 51.7 & 68.2\\
		\midrule
		\multicolumn{8}{l}{Panel B: market-wide P/E}\\
		2002--2012 & 24.8 & 17.8 & 24.5 & 30.2 & 45.5 & 34.4 & 66.9\\
		2002--2007 & 26.8 & 18.2 & 30.0 & 29.3 & 50.9 & 33.1 & 75.1\\
		2008--2012 & 21.7 & 17.5 & 20.1 & 31.9 & 39.5 & 38.5 & 55.1\\
		\bottomrule
	\end{tabular}
	\par\vspace{5pt}
	\small U = underpriced，O = overpriced。即便被指标判定为 underpriced，高 LA 股票的 P/E 也明显偏高，这与 Table 11--12 的解释一致：高 LA 组整体处于偏贵的位置。
\end{frame}

\section{共同启示}

\begin{frame}{两篇文章共同给出的研究路径}
	把两篇文章合起来看，可以提炼出一个相对完整的研究步骤：
	\begin{enumerate}
		\item 从已经被国际文献确认的经典异象出发——MAX 和 IVOL 在美国和欧洲都已有大量证据。
		\item 找到中国市场与发达市场最不一样的地方：散户主导、涨跌停、卖空限制、融资融券标的、股指期货覆盖范围。
		\item 把这些制度差异转化为可量化的变量：例如 Nartea 对 MAX 做涨跌停适配，Gu 用六个组件构造 LA 指数。
		\item 提出的问题不止是 ``中国是否也存在''，而是 ``机制在这里有没有发生变化''——MAX 会不会吸收 IVOL？LA 会不会放大 IVOL？
		\item 实证按照单变量排序、双/三重排序、Fama--MacBeth、替代定义、机制排除的顺序逐层推进。
	\end{enumerate}
\end{frame}

\begin{frame}{对写"中国 anomaly 文章"的几点提醒}
	\begin{itemize}
		\item 单纯把一个美国异象搬到中国跑一遍，贡献往往很有限；关键是中国的制度或投资者结构有没有帮助识别背后的机制。
		\item 变量不在多，而在能不能给出一个清晰的、可证伪的预测。例如 Gu et al.\ 的 ``高 LA 组 IVOL premium 应当更强''，就是一个 sharp prediction。
		\item 表格之间要有递进关系：先看现象是否存在，再看是否被现有特征吸收，再做机制拆分，最后做稳健性，而不是反复堆显著性。
		\item 当一个机制解释被否定时，需要进一步说明替代机制为什么不够，而不是停在 ``加上控制变量后依然显著''。
	\end{itemize}
\end{frame}

\begin{frame}[c]{两篇文章的差别在哪里}
	\tighttable
	\begin{tabularx}{\textwidth}{lYY}
		\toprule
		维度 & Nartea et al. & Gu et al.\\
		\midrule
		文章重心 & 验证 MAX 效应在中国是否存在，并讨论它和 IVOL 的关系。 & 解释 IVOL premium 为什么在中国更强：套利限制放大了 mispricing。\\
		机制变量 & MAX 本身就是 lottery-like payoff 的代理，背景机制是中国投资者的风险偏好。 & LA 指数是核心机制变量，直接进入排序和交互回归。\\
		最具说服力的表 & Table 1、3、5、9、10。 & Table 3、6、7、9、10、11、12。\\
		结论的口径 & MAX 与 IVOL 在中国是并存的两个异象。 & 高 LA 让 IVOL 异象更强、更持久；已有的几种解释都不够充分。\\
		\bottomrule
	\end{tabularx}
\end{frame}

\begin{frame}{课堂讨论的三个落点}
	\begin{enumerate}
		\item \textbf{MAX 和 IVOL 的高相关性必须正视。} 不互相控制，就无法判断驱动横截面定价的到底是彩票偏好还是特质波动。
		\item \textbf{套利限制是中国市场最有价值的制度特征之一。} 它能把 ``异常收益'' 的问题转化为 ``价格为什么不能被及时纠正'' 的机制问题。
		\item \textbf{在中国做异象研究，不能只是 replication。} 最值得记住的部分恰恰是与美国证据不同的地方——MAX 没能吸收 IVOL；IVOL 在高 LA 组里更强、更持久。
	\end{enumerate}
\end{frame}

\begin{frame}{小结}
	\begin{beamercolorbox}[sep=0.6em,wd=\linewidth,rounded=true]{postit}
		把两篇文章放在一起，可以看到一种中国资产定价的实证范式：借助中国投资者结构和交易制度摩擦，去重新审视经典资产定价异象的机制边界。
	\end{beamercolorbox}
	\vspace{8pt}
	\begin{itemize}
		\item Nartea et al.：极端收益越高，未来收益越低；彩票偏好解释有支持，但 MAX 与 IVOL 在中国是两个并存的异象。
		\item Gu et al.：高 IVOL 股票未来收益低，并且在套利越受限的股票中，这一负收益更强、更持久。
		\item 方法上最值得学习的，是它们打磨表格链的方式：从 ``异象是否存在'' 一步步走到 ``机制如何识别''。
	\end{itemize}
\end{frame}

\end{document}
